\documentclass[12pt]{article}
\usepackage[utf8]{inputenc}
\usepackage[T1]{fontenc}
\usepackage{hyperref}
\title{Tifal Language Resources for K-12 Curriculum Development}
\author{Compiled by Claude Sonnet 5 (claude-sonnet-5) for Andr\'as Kornai}
\date{}
\begin{document}
\maketitle
\section{Introduction}
Tifal (ISO 639-3 tif) is a Central Mountain Ok language of Papua New Guinea. It is spoken in the Telefomin district (Sandaun) and in North Fly district (Western Province). It is closely related to Telefol. Speaker counts differ between sources: about 3,600, 4,000 (2003) and 6,900 (Ethnologue 2016, via Joshua Project). Dialects listed are Tifal (Tifalmin), Urap (Urapmin) and Atbal (Atbalmin).
\section{Primary Educational Resources}
None found.
\section{Linguistic and Reference Materials}
\begin{itemize}
\item Tifal numeral system, a 27-unit body-part tally: \url{https://lingweb.eva.mpg.de/channumerals/Tifal.htm}
\item Mountain Ok family overview: \url{https://newguineaworld.linguistik.uzh.ch/families/trans-new-guinea/central-west-new-guinea/digul-river-ok/ok/mountain-ok/start}
\item Telefol language (a close relative, useful for comparison): \url{https://en.wikipedia.org/wiki/Telefol_language}
\end{itemize}
\section{Religious and Community Resources}
Joshua Project profile: \url{https://joshuaproject.net/people_groups/15479/PP}.
\section{Digital and Audio Resources}
None found.
\section{Curriculum Development Context}
No Tifal teaching materials were found. Related Ok languages such as Telefol may be a useful starting point for orthography and materials design, but this was not checked.
\end{document}